\lecture{2}{Wed 10 Jan 2024 13:31}{}

On a normed linear space, the algebraic operations (addition and scalar multiplication) are continuous, and \textit{translations are isometries}: $U(x, r) = x + U(0, r)$.

\begin{remark}
	Different norms generate different ``geometries'' in $\mathbb{R}^n$.  However, they give the same topology.

	You can nest all the different balls in each other, so topologies are the same.
\end{remark}

\begin{example}
	 We can define different norms on the linear space $X = \{f: \mathbb{N} \to  \mathbb{C}, f \text{ finite support}\} $.
	 \begin{enumerate}
	 	\item $\|f\|_2 = \left( \sum_{n \in \mathbb{N}} |f(n)|^2 \right) ^{1 / 2}$
		\item $\| f\|_1 = \sum_{n \in \mathbb{N}} |f(n)|$
		\item $\| f\|_p = \left( \sum _{n \in \mathbb{N}} |f(n)|^p \right) ^{1 / p}$
		\item $\| f\|_{\infty } = \sup _n |f(n)|$
	 \end{enumerate}
	 Also have continuous versions $\|f\|_p = \left( \int _0^\infty  |f(t)|^p \right) ^{1 / p}$
\end{example}

Now we want to generalize. 

\begin{definition}[Linear Space associated to a set]
	\label{defn:LinearSpaceAssociatedToASet}
	Linear space associated to a set $T$ is $\mathbb{C} [ T] = \{f: T \to \mathbb{C}: \text{supp} f \text{ finite}\} $
	
	A basis consists of Dirac functions  $\delta_t, t \in T$.

	Consider counting measure $\mu$ on $T$, each point has measure $1$. Now make $\mathbb{C}[T]$ into a normed space by setting
	\[
		\|f\|_p = \left( \int _T |f(t)|^p d \mu \right) ^{1 / p} 
	.\] 
\end{definition}

\begin{definition}[Banach Space]
	\label{defn:BanachSpace}
	A \textit{Banach space} is a normed linear space which is complete with respect to the metric $d(x, y) = \|x - y\|$.
\end{definition}

\begin{remark}
	Note \textit{homogeneity} of a normed metric space. Metric structure looks the same around all points. All translations are isometric homeomorphisms:
	\[
		U(x, r) = x + U(0, r)
	.\] 
	\textbf{Homothety Property}: $\lambda \| x\| = \| \lambda x\|$
\end{remark}

These properties are crucial for understanding linear operators and functionals. 

\begin{example}
	\begin{enumerate}
		\item If $T$ is a compact Hausdorff space, 
			\[
				C(T) = \{f: T \to \mathbb{C}, f \text{ continuous}\} 
			\] 
			is a Banach space wrt the sup norm. \textbf{Fundamental example for spectral theory.}
		\item For $T$ nonempty set, (possibly infinite),
			\[
				l^{\infty } (T, \mathbb{C}) = l^{\infty }(T) = \{f: T \to  \mathbb{C} , f \text{ bounded}\} 
			\] 
			is a Banach space with sup norm. 

			For $T$ compact Hausdorff, $C(T) \subset l^{\infty }(T)$ with induced norm.
		\item If $A$ is a Banach space, then
			\[
				l^{\infty }(T, A) = \{ f: T \to  A, f \text{ bounded}\} 
			\] 
			is a Banach space with the sup norm.
	\end{enumerate}
\end{example}

\begin{proposition}[Functional Space is Complete]
	\label{prop:FunctionalSpaceIsComplete}
	$l^\infty (T, A)$ is complete.
\end{proposition}
\begin{proof}
	Suppose $f_n$ is Cauchy, find $n_0$ such that $\|f_n - f_m\|_\infty  < \varepsilon $ for $n, m \ge  n_0(\varepsilon)$. For all $t \in  T$, $f_n(t)$ is a Cauchy sequence in $A$ since $\|f_n(t) - f_m(t)\|_A \le  \|f_n - f_m\|_\infty < \varepsilon $, and hence convergent to some $f(t) \in A$. 

	We constructed a function $f: T \to A$ such that $f_n \to  f$ point wise. Now we want uniform convergence. 
	Take $m \to  \infty $, we get $\| f_n(t) - f(t)\|_A < \varepsilon$ for all $n \ge  n_0, t \in T$. Hence for $n > n_0$,
	$\|f_n - f\|_\infty = \sup _{t \in T} \|f_n(t) - f(t)\|_A \le \varepsilon
	.$ 
\end{proof}

\begin{proposition}
	\label{prop:InducedCompleteness}
	If $Y \subset X$  is closed in $X$ and $X$ complete, then $Y$ complete.
\end{proposition}
\begin{proof}
	$X$ complete, so any Cauchy sequence in $Y$ converges in $X$. $Y$ closed, so the limit is inside $Y$.
\end{proof}

\begin{definition}[Isometry]
	\label{defn:Isometry}
	An \textit{isometry} is a map $f: X \to  Y$ that preserves metric.
\end{definition}

\begin{definition}[Completion]
	\label{defn:Completion}
	A \textit{completion} of a metric space $X$ is a complete metric space $Y$ together with an isometric embedding $j: X \to  Y$ such that $j(X)$ is dense in $Y$. 
\end{definition}

\begin{lemma}[Uniqueness of Completion]
	\label{lem:UniquenessOfCompletion}
	Completions of metric spaces are unique up to isometry.
\end{lemma}
\begin{proof}
% https://quiver.local:8080/#q=WzAsNSxbMiwwLCJqXzEoWCkiXSxbMCwwLCJYIl0sWzIsMiwial8yKFgpIl0sWzMsMCwiWV8xIl0sWzMsMiwiWV8yIl0sWzEsMCwial8xIl0sWzEsMiwial8yIiwyXSxbMCwyLCJqXzJqXzFeey0xfSIsMl0sWzMsNCwiRiIsMCx7ImNvbG91ciI6WzAsNjAsNjBdLCJzdHlsZSI6eyJib2R5Ijp7Im5hbWUiOiJkYXNoZWQifX19LFswLDYwLDYwLDFdXSxbMCwzLCIiLDEseyJjb2xvdXIiOlswLDYwLDYwXSwic3R5bGUiOnsidGFpbCI6eyJuYW1lIjoiaG9vayIsInNpZGUiOiJ0b3AifX19XSxbMiw0LCIiLDEseyJzdHlsZSI6eyJ0YWlsIjp7Im5hbWUiOiJob29rIiwic2lkZSI6InRvcCJ9fX1dLFswLDQsIiIsMCx7ImNvbG91ciI6WzAsNjAsNjBdfV1d
\[\begin{tikzcd}
	X && {j_1(X)} & {Y_1} \\
	\\
	&& {j_2(X)} & {Y_2}
	\arrow["{j_1}", from=1-1, to=1-3]
	\arrow["{j_2}"', from=1-1, to=3-3]
	\arrow["{j_2j_1^{-1}}"', from=1-3, to=3-3]
	\arrow["F", color={rgb,255:red,214;green,92;blue,92}, dashed, from=1-4, to=3-4]
	\arrow[color={rgb,255:red,214;green,92;blue,92}, hook, from=1-3, to=1-4]
	\arrow[hook, from=3-3, to=3-4]
	\arrow[color={rgb,255:red,214;green,92;blue,92}, from=1-3, to=3-4]
\end{tikzcd}\]

Because
\begin{itemize}
	\item $j_1(X)$ admits an isometric embedding to $Y_2$, by going around the diagram 
	\item $j_1(X)$ is dense in $Y_1$
\end{itemize}
Hence, by the fundamental result, there exists an isometric embedding $Y_1 \to Y_2$.

Do this oppositely, and we see that $Y_2 \to Y_1$ is also an isometric embedding, and those two embeddings are inverse to each other. Hence $Y_1 $ and $Y_2$ are same up to isometry.

One can also directly extend the bijective isometric embedding $j_2 j_1^{-1}$ to $F$.

\end{proof}

\begin{proposition}
	\label{prop:EmbeddingMetricSpace}
	Every metric space $(X, d)$ embeds isometrically as a subspace of $l^{\infty }(T)$ for some set $T$.
	\[
		\exists  j: X \to  l^\infty (T), \|j(y) - j(y')\|_\infty  = d(y, y')
	.\] 
\end{proposition}
\begin{proof}
	Define $j: X \to  l^\infty (X, \mathbb{R})$ as follows: fix $x_0 \in X$, then 
	\[
		j: x \mapsto j(x) = f_x = d(\cdot , x) - d(\cdot , x_0) 
	.\] 
	To verify $f_x \in l^\infty (X, \mathbb{R})$, note that
	 \[
		\|f_x\|_{\infty } = \sup |d(\cdot , x) - d(\cdot , x_0)|  = d(x, x_0) < \infty 
	.\] 
	The last equality is due to triangle inequality.

	$j$ is clearly injective. To verify $j$ is an isometry,
	\[
		\|j x - j y\|_{\infty } = \|d(\cdot , x) - d(\cdot , x_0) - d(\cdot , y) + d(\cdot , x_0)\|_{\infty } = d(x, y)
	.\] 
	The last equality is again due to triangle inequality.
\end{proof}

\begin{corollary}
	Every metric space $(Y, d)$ admits a completion $(X, d)$ which is unique up to an isometry.
\end{corollary}
\begin{proof}
	We already know the completions are unique (by Lemma~\ref{lem:UniquenessOfCompletion} ). We only need to demonstrate that such a completion exists. But we also know from the last proposition, that $(X, d)$ is isometrically embedded into $l^\infty (X)$ which is complete. Hence, by Proposition~\ref{prop:InducedCompleteness}, the closure of its embedded image is a complete subspace of $l^\infty (X)$, which is a completion of $(X, d)$:
	\[
		(X, d(x, x')) = \left( \overline{j(Y)} , \|x - x'\|_\infty  \right) 
	.\] 
\end{proof}

\begin{definition}
	\label{defn:CompletionOfBanachSpace}
	
	A \textit{completion} of a normed linear space $X$ is a Banach space $Y $ with an \textbf{isometric} \textbf{linear} embedding $j: X \to  Y$ such that $\overline{j(X)}  = Y$.
\end{definition}

\begin{remark}
	The completions of normed spaces from last example are
\[
\begin{aligned}
& \ell^p(\mathbb{N})=\left\{f: \mathbb{N} \rightarrow \mathbb{C}:\|f\|_p=\left(\sum_{n \in \mathbb{N}}|f(n)|^p\right)^{1 / p}<\infty\right\} . \\
& \ell^1(\mathbb{N})=\left\{f: \mathbb{N} \rightarrow \mathbb{C}:\|f\|_1=\sum_{n \in \mathbb{N}}|f(n)|<\infty\right\} \\
& C_0(\mathbb{N})=\left\{f: \mathbb{N} \rightarrow \mathbb{C}: \lim _{n \rightarrow \infty} f(n)=0\right\} .
\end{aligned}
\]
\end{remark}

\begin{example}(Lebesgue measurable functions)
	Consider $T = [0,1]$ with Lebesgue measure $\lambda$. $L^1[0,1]$ is equivalence classes of all Lebesgue measurable functions that are  $L^1$ finite.

	Let $(\mathscr{S}, \|\cdot \|_1)$ be normed linear space consisting of simple functions of the form $f = \sum _{i = 1}^n a_i \chi_{T_i}$, $T_i$ Lebesgue measurable. Then $L^1[0,1]$ is the Banach space completion of $(\mathscr{S}, \|\cdot \|_1)$.
\end{example}
A note from measure theory: this construction precisely gives the measurable morphisms from $([0,1], \lambda)$ to $(\mathbb{C}, \mu_{\text{Borel}})$ with finite $L^1$ norm. 
\todo{Check if this statement is precise.}

\begin{example}
	To illustrate the extension theorem, we define the integral for a continuous functions $f: [0,1] \to A$ a Banach space.

	Consider $\mathscr{S}$ be linear subspace of $l^\infty (T, A)$ consisting of simple functions. Define
	\[
		\phi(f) = \sum_i a_i \lambda(T_i)
	.\] 
	where $\lambda$ is Lebesgue measure. where $\lambda$ is the Lebesgue measure. Check that $\varphi$ is linear and contractive.
\[
\begin{gathered}
\|\varphi(f)\| \leqslant \sum_{l \in I}\left\|a_i\right\| \lambda\left(T_i\right) \leqslant \sup _{i \in I}\left\|a_i\right\|=\|f\|, \quad \text { hence } \\
\|\varphi(f)-\varphi(g)\| \leqslant\|f-g\| .
\end{gathered}
\]

By fundamental Theroem~\ref{thm:FundamentalResult}$, \varphi$ it extends to $\Phi: \overline{\mathcal{S}} \rightarrow A$.
$\overline{\mathcal{S}}$ contains all Borel functions $g: T \rightarrow A$ with relatively compact image (why?) and in particular
\[
C(T, A)=\{f: T=[0,1] \rightarrow A, \text { continuous }\} \subseteq \overline{\mathcal{S}}
\]

If $g \in \overline{\mathcal{S}}$ define
\[
\int_T g d \lambda=\Phi(g)
\]

Terminology: $g$ Borel functions means that $g^{-1}(U)$ is a Borel set for each open subset $U$ of $A$.
\end{example}

\todo{I currently do not know how to show the uniform limit of simple functions have relative compact image. How to do that for $\mathbb{R}$-valued functions?}


